% Work/MAT | 1. rocnik | nova hodina 31: Castecne odmocnovani.
% Nova posloupnost: 29 druha/treti odmocnina, 30 racionalni exponent,
% 31 castecnne odmocnovani, 32 usmernovani jednoclenneho jmenovatele,
% 33 usmernovani sdruzenym vyrazem, 34 slozene vyrazy, 35 vedecky zapis.
% Ucebnice: ucebnice/01_Zakladni_poznatky/Mocniny_a_odmocniny/Mocniny a Odmocniny.tex
% (Odmocniny, Castecne odmocnovani, Caste chyby a omyly).
% Metodika: Realisticky.cz 010207 Druha odmocnina, 020712 Pocitani s odmocninami I.
% Vsechna zadani nove sestavena, z nejsou doslovne prepisy zdroju.
\documentclass[11pt,a4paper]{article}
% Spolecny styl pro ucitelske pripravy MAT (Work/MAT)
% Verze 1.1, 2026-10-23
\NeedsTeXFormat{LaTeX2e}
\ProvidesPackage{priprava}[2026/10/23 v1.1 Sablona ucitelskych priprav MAT]

% Zakladni typografie a cestina
\RequirePackage[T1]{fontenc}
\RequirePackage[utf8]{inputenc}
\RequirePackage[czech]{babel}
\RequirePackage{lmodern}
\RequirePackage{microtype}
\RequirePackage{amsmath,amssymb}
\RequirePackage{xcolor}
\RequirePackage{enumitem}
\RequirePackage[most]{tcolorbox}
\RequirePackage{tabularx,booktabs,array}
\RequirePackage{fancyhdr}
\RequirePackage{needspace}
\RequirePackage{geometry}
\RequirePackage[hidelinks]{hyperref}

\geometry{a4paper,top=17mm,bottom=18mm,left=23mm,right=23mm}
\definecolor{notegray}{gray}{0.39}
\setlength{\parindent}{0pt}
\setlength{\parskip}{2.5pt plus 1pt minus 0.5pt}
\setlist[itemize]{topsep=2pt,itemsep=1pt,parsep=0pt,leftmargin=1.45em}
\setlist[enumerate]{topsep=2pt,itemsep=2pt,parsep=0pt,leftmargin=1.65em}
\setcounter{secnumdepth}{0}
\renewcommand{\arraystretch}{1.12}

% Zahlavi: vlevo rocnik | globalni tema, vpravo cislo: nazev hodiny
% Zapati zustava zcela prazdne.
\pagestyle{fancy}
\fancyhf{}
\renewcommand{\headrulewidth}{0.35pt}
\renewcommand{\footrulewidth}{0pt}
\setlength{\headheight}{14pt}
\newcommand{\nastavPripravu}[4]{%
  \fancyhead[L]{\small #1\enspace\textbar\enspace #2}%
  \fancyhead[R]{\small #3: #4}%
}

% Minimalisticke cernobile boxy. Kazda definice zvlast.
\tcbset{commonlesson/.style={enhanced,breakable=false,colback=white,colframe=black,
  colbacktitle=white,coltitle=black,boxrule=0.45pt,
  arc=0.8mm,outer arc=0.8mm,left=3mm,right=3mm,top=1.6mm,bottom=1.6mm,
  toptitle=1mm,bottomtitle=0.7mm,fonttitle=\bfseries,
  attach boxed title to top left={xshift=2.5mm,yshift=-1.15mm},
  boxed title style={colback=white,colframe=white,boxrule=0pt,left=1mm,right=1mm}}}
\newtcolorbox{zadaniBox}[1][]{commonlesson,title={Zadání},#1}
\newtcolorbox{prikladBox}[1][]{commonlesson,title={Příklad},#1}
\newtcolorbox{definiceBox}[1][]{commonlesson,title={Definice},#1}

% Spolecny box pro zadani a jeho reseni, bez mezery mezi castmi.
% Pouziti:
%   \begin{prikladReseniBox}
%     Zadani ulohy.
%     \tcblower
%     \begin{align*} ... \end{align*}
%   \end{prikladReseniBox}
% Spodni cast ma jemny sedy podklad a delici linku; lze v ni sazet
% vice radku matematiky. Obe casti maji jen jeden spolecny ramecek.
\tcbset{solutionlesson/.style={
  commonlesson,
  enhanced jigsaw,
  breakable,
  colbacklower=black!3,
  segmentation style={solid,draw=black!25,line width=0.35pt},
  fontlower=\small\color{notegray},
  before lower={\noindent\textbf{Řešení:}\enspace},
  leftlower=3mm,rightlower=3mm
}}
\newtcolorbox{prikladReseniBox}[1][]{solutionlesson,title={Příklad},#1}
\newtcolorbox{zadaniReseniBox}[1][]{solutionlesson,title={Zadání},#1}

% Pomucky pro autora pripravy: sedive poznamky; \reseni pro starsi pripravy
\newcommand{\poznamka}[1]{\par\begingroup\small\color{notegray}\itshape
  \textnormal{\textbf{Poznámka pro učitele.}}\enspace #1\par\endgroup}
\newcommand{\reseni}[1]{\par\begingroup\small\color{notegray}
  \textbf{Očekávané řešení:}\enspace #1\par\endgroup}
\newcommand{\blok}[1]{\needspace{5\baselineskip}\subsection*{#1}}
\newcommand{\vyrok}[1]{\uv{#1}}

\nastavPripravu{1. ročník}{Mocniny a odmocniny}{31}{Částečné odmocňování}
\begin{document}
\poznamka{\textbf{45minutové jádro:} 4 min opakování; 8 min odvození a Příklad 1;
9 min Samostatná práce 1 a)--f); 7 min Příklad 2 a Samostatná práce 2 a)--d);
8 min sčítání podobných odmocnin a Příklad 3; 6 min Samostatná práce 3 a)--c);
3 min shrnutí a zadání domácí práce. Příklad 4, rozšíření, další položky a delší domácí úlohy
jsou zásobník. Tabulka pravidel z 30. hodiny se zde \textbf{nemění}: používáme ji.}

\blok{Proč vůbec odmocňujeme jen částečně?}
Z hodiny 29 známe $\sqrt{36}=6$, ale $\sqrt{72}$ není celé číslo.
Hledáme proto uvnitř \emph{co největší vhodný čtverec} a rozložíme radikand:
\[
\sqrt{72}=\sqrt{36\cdot2}=\sqrt{36}\sqrt2=\boxed{6\sqrt2}.
\]
\textbf{Výhoda:} přesná hodnota zůstane zachována bez zbytečného zaokrouhlování.
U třetí odmocniny hledáme místo čtverců \emph{třetí mocniny}.

\begin{definiceBox}[title={Pravidlo: odmocnina ze součinu}]
Pro $a,b\ge0$ platí $\sqrt{ab}=\sqrt a\,\sqrt b$.
Pro libovolná $a,b\in\mathbb R$ platí
$\sqrt[3]{ab}=\sqrt[3]a\,\sqrt[3]b$.
\textbf{Částečně odmocnit} znamená vhodný čtverec nebo třetí mocninu
vyjmout před odmocnítko; u druhé odmocniny musí být vyjmutá hodnota nezáporná.
\end{definiceBox}
\textbf{Odvození pro druhou odmocninu:}
Pro $a,b\ge0$ je $\sqrt a\sqrt b\ge0$ a
$(\sqrt a\sqrt b)^2=ab$. Podle definice druhé odmocniny tedy
$\sqrt a\sqrt b=\sqrt{ab}$.
\poznamka{U třetí odmocniny funguje stejná úvaha s třetí mocninou;
na rozdíl od druhé není nutná nezápornost radikandu. Žáci mají vyslovit,
\emph{jaký násobek dokonalého čtverce / krychle hledají}.}

\begin{prikladReseniBox}[breakable=false,title={Příklad 1: vyjmutí čtverce}]
Částečně odmocni $\sqrt{48}$, $\sqrt{108}$, $\sqrt{245}$.
\tcblower
\[
\sqrt{48}=\sqrt{16\cdot3}=4\sqrt3,\quad
\sqrt{108}=\sqrt{36\cdot3}=6\sqrt3,\quad
\sqrt{245}=\sqrt{49\cdot5}=7\sqrt5.
\]
Pokud by pod odmocnítkem zůstal další čtvercový činitel, lze pokračovat.
\end{prikladReseniBox}

\begin{zadaniReseniBox}[breakable=false,title={Samostatná práce 1: druhé odmocniny}]
Částečně odmocni; nepoužívej kalkulačku.
\begin{enumerate}[label=\alph*)]
\item $\sqrt{12}$; \item $\sqrt{18}$; \item $\sqrt{50}$; \item $\sqrt{75}$;
\item $\sqrt{128}$; \item $\sqrt{300}$; \item $\sqrt{800}$; \item $\sqrt{0{,}72}$.
\end{enumerate}
\tcblower
\begin{enumerate}[label=\alph*)]
\item $2\sqrt3$; \item $3\sqrt2$; \item $5\sqrt2$; \item $5\sqrt3$;
\item $8\sqrt2$; \item $10\sqrt3$; \item $20\sqrt2$;
\item $\sqrt{0{,}36\cdot2}=0{,}6\sqrt2$.
\end{enumerate}
\end{zadaniReseniBox}

\blok{Třetí odmocnina: hledáme dokonalou krychli}
\begin{prikladReseniBox}[breakable=false,title={Příklad 2: vyjmutí třetí mocniny}]
Částečně odmocni $\sqrt[3]{54}$, $\sqrt[3]{-128}$ a $\sqrt[3]{250}$.
\tcblower
\[
\begin{aligned}
\sqrt[3]{54}&=\sqrt[3]{27\cdot2}=3\sqrt[3]2,\\
\sqrt[3]{-128}&=\sqrt[3]{(-64)\cdot2}=-4\sqrt[3]2,\\
\sqrt[3]{250}&=\sqrt[3]{125\cdot2}=5\sqrt[3]2.
\end{aligned}
\]
U záporného čísla třetí odmocnina existuje, protože jde o \emph{lichou} odmocninu.
\end{prikladReseniBox}

\begin{zadaniReseniBox}[breakable=false,title={Samostatná práce 2: třetí odmocniny}]
\begin{enumerate}[label=\alph*)]
\item $\sqrt[3]{16}$; \item $\sqrt[3]{81}$;
\item $\sqrt[3]{-54}$; \item $\sqrt[3]{375}$;
\item $\sqrt[3]{-500}$; \item $\sqrt[3]{432}$.
\end{enumerate}
\tcblower
\begin{enumerate}[label=\alph*)]
\item $2\sqrt[3]2$; \item $3\sqrt[3]3$;
\item $-3\sqrt[3]2$; \item $5\sqrt[3]3$;
\item $-5\sqrt[3]4$; \item $6\sqrt[3]2$.
\end{enumerate}
\end{zadaniReseniBox}

\blok{Sčítání podobných odmocnin}
$3\sqrt2+5\sqrt2=8\sqrt2$ funguje podobně jako $3x+5x=8x$.
\textbf{Sčítat koeficienty smíme až tehdy}, když mají odmocniny stejný index i radikand.
\begin{prikladReseniBox}[breakable=false,title={Příklad 3: nejprve odmocni částečně, pak slučuj}]
Uprav $\sqrt{72}-\sqrt{18}+2\sqrt8$.
\tcblower
$\sqrt{72}=6\sqrt2$, $\sqrt{18}=3\sqrt2$, $2\sqrt8=4\sqrt2$.
Proto $(6-3+4)\sqrt2=\boxed{7\sqrt2}$.
\end{prikladReseniBox}
\begin{zadaniReseniBox}[breakable=false,title={Samostatná práce 3: slučuj pouze stejné odmocniny}]
\begin{enumerate}[label=\alph*)]
\item $\sqrt{27}+\sqrt{12}$;
\item $2\sqrt{50}-3\sqrt8$;
\item $\sqrt{75}-\sqrt{48}+\sqrt3$;
\item $\sqrt[3]{16}+\sqrt[3]{54}-\sqrt[3]2$;
\item $\sqrt2+\sqrt3$ (lze sloučit?).
\end{enumerate}
\tcblower
\begin{enumerate}[label=\alph*)]
\item $3\sqrt3+2\sqrt3=5\sqrt3$;
\item $10\sqrt2-6\sqrt2=4\sqrt2$;
\item $5\sqrt3-4\sqrt3+\sqrt3=2\sqrt3$;
\item $2\sqrt[3]2+3\sqrt[3]2-\sqrt[3]2=4\sqrt[3]2$;
\item Ne, jde o různé radikandy. Zejména $\sqrt2+\sqrt3\ne\sqrt5$.
\end{enumerate}
\end{zadaniReseniBox}

\blok{Když je pod odmocninou proměnná}
\begin{prikladReseniBox}[breakable=false,title={Příklad 4: rozdíl mezi druhou a třetí odmocninou}]
Pro reálné $x$ uprav $\sqrt{18x^2}$ a $\sqrt[3]{-54x^3}$.
\tcblower
$\sqrt{18x^2}=3\sqrt2\,\sqrt{x^2}=\boxed{3|x|\sqrt2}$.
Dále $\sqrt[3]{-54x^3}=-3x\sqrt[3]2$.
U druhé odmocniny \emph{nelze obecně} místo $|x|$ napsat $x$.
\end{prikladReseniBox}
\poznamka{Příklad 4 je \textbf{rezerva pro rychlejší žáky}. Ověřte třeba $x=-2$.
Připomeň $\sqrt{x^2}=|x|$. Jde o metodicky zásadní pravidlo, ne o formalitu.}

\begin{zadaniReseniBox}[breakable=false,title={Rozšiřující příklad 1: součiny a proměnné}]
\begin{enumerate}[label=\alph*)]
\item $\sqrt{72a^2}$ pro $a\in\mathbb R$;
\item $\sqrt{27x^4y^2}$ pro $x,y\in\mathbb R$;
\item $\sqrt[3]{-250t^6}$ pro $t\in\mathbb R$;
\item $\sqrt{2}\sqrt{8}+\sqrt{3}\sqrt{12}$.
\end{enumerate}
\tcblower
\begin{enumerate}[label=\alph*)]
\item $6|a|\sqrt2$;
\item $3x^2|y|\sqrt3$;
\item $-5t^2\sqrt[3]2$;
\item $4+6=10$.
\end{enumerate}
\end{zadaniReseniBox}

\blok{Domácí procvičování}
\begin{zadaniReseniBox}[breakable=false,title={Domácí úkol 1: částečné odmocňování}]
\begin{enumerate}[label=\alph*)]
\item $\sqrt{200}$; \item $\sqrt{147}$; \item $\sqrt[3]{-192}$;
\item $\sqrt[3]{135}$; \item $3\sqrt{12}-\sqrt{75}$;
\item $2\sqrt{45}+\sqrt{20}-\sqrt5$.
\end{enumerate}
\tcblower
\begin{enumerate}[label=\alph*)]
\item $10\sqrt2$; \item $7\sqrt3$; \item $-4\sqrt[3]3$;
\item $3\sqrt[3]5$; \item $6\sqrt3-5\sqrt3=\sqrt3$;
\item $6\sqrt5+2\sqrt5-\sqrt5=7\sqrt5$.
\end{enumerate}
\end{zadaniReseniBox}
\poznamka{\textbf{Kontrola na konci:} Proč je $\sqrt{50}=5\sqrt2$ a ne $25\sqrt2$?
Proč $\sqrt{a^2}=|a|$? V následující hodině využijeme částečné
odmocňování před usměrňováním zlomků.}
\end{document}
